\chapter{Graphs and How We Represent Them}

A graph is about as spare an object as computing has: a set of things and
a set of relationships between them. In the preface we watched how a traversal
explores outward along connected vertices. Before we can analyze or execute any
such algorithm, however, we must decide how vertices and relationships are defined
mathematically and how they live inside computer memory. A bad representation can
turn a linear-time algorithm into a quadratic one without changing a single
line of its logic. This opening chapter establishes our core vocabulary, compares
the standard data structures for graphs, and fixes the asymptotic cost model that
governs every running time in this book.

\chapterline

\section{What a graph is}
\label{ch01:formal}

\begin{intuitionbox}
Draw some dots. Connect some of them with lines. You now have a graph. The dots
are \emph{vertices} (or nodes), the lines are \emph{edges}. Everything else ---
direction, weight, the prohibition or tolerance of repeated edges --- is a
decision you make about what the dots and lines are allowed to mean. The art of
graph modeling is choosing those meanings so that the question you care about
becomes a standard graph question with a known answer.
\end{intuitionbox}

\begin{definition}[Undirected graph]
\label{ch01:def-undirected}
An \emph{undirected graph} is a pair $\Gph = (V, E)$ where $V$ is a finite set
of \emph{vertices} and $E$ is a set of \emph{edges}, each edge being a
two-element subset $\set{u,v} \subseteq V$ with $u \neq v$. We write
$n = \abs{V}$ and $m = \abs{E}$.
\end{definition}

The requirement $u \neq v$ rules out \emph{self-loops} (an edge from a vertex to
itself), and using a \emph{set} of edges rules out \emph{parallel edges} (two
distinct edges joining the same pair). A graph with neither is called
\emph{simple}. When we want loops or parallel edges back we will say
\emph{multigraph}; for now, simple is the default, and unless a chapter says
otherwise, ``graph'' means ``simple graph.''

\begin{definition}[Directed graph]
\label{ch01:def-directed}
A \emph{directed graph} (or \emph{digraph}) is a pair $\Gph = (V, E)$ where
$E \subseteq V \times V$ is a set of ordered pairs. An edge $(u,v)$ points
\emph{from} its \emph{tail} $u$ \emph{to} its \emph{head} $v$. We draw it as an
arrow $u \to v$.
\end{definition}

In a digraph the pair is ordered, so $(u,v)$ and $(v,u)$ are different edges and
both may be present. The two pictures below show the distinction. On the left an
undirected edge between $a$ and $b$; on the right two directed edges, one each
way, which together carry the same connectivity but twice the bookkeeping.

\begin{center}
\begin{tikzpicture}[every node/.style={circle,draw=nodeblue,fill=blue!8,inner sep=2pt,minimum size=20pt,font=\small},
  >={Stealth[length=5pt]}]
  \node (a) at (0,0) {$a$};
  \node (b) at (2,0) {$b$};
  \draw[edgegray,thick] (a) -- (b);
  \node[draw=none,fill=none] at (1,-0.9) {undirected $\set{a,b}$};

  \node (c) at (5,0) {$a$};
  \node (d) at (7,0) {$b$};
  \draw[->,edgegray,thick] (c) to[bend left=22] (d);
  \draw[->,edgegray,thick] (d) to[bend left=22] (c);
  \node[draw=none,fill=none] at (6,-0.9) {directed $(a,b),(b,a)$};
\end{tikzpicture}
\end{center}

\begin{definition}[Weighted graph]
\label{ch01:def-weighted}
A \emph{weighted graph} adds a function $w : E \to \R$ (or to $\R_{\ge 0}$, or to
$\Z$, depending on the setting) assigning each edge a number. We write $w(u,v)$
or $w_e$. Weights model distance, cost, capacity, similarity, resistance ---
whatever the application measures along an edge.
\end{definition}

Weight is orthogonal to direction: there are weighted undirected graphs (road
networks with mileage) and weighted directed graphs (currency exchange rates).
Whether weights may be negative matters enormously --- Dijkstra needs them
nonnegative, Bellman--Ford tolerates negatives but not negative cycles --- and
we will be careful about that in the shortest-path chapters.

\begin{definition}[Multigraph]
\label{ch01:def-multigraph}
A \emph{multigraph} replaces the edge \emph{set} by an edge \emph{multiset}, or
equivalently by a set $E$ of edge identities together with an
\emph{endpoint map} $\partial$ sending each edge to its (un)ordered pair of
endpoints. This permits parallel edges and, if we also drop the $u \neq v$
restriction, self-loops.
\end{definition}

The endpoint-map formulation is worth a second look because it is the clean way
to talk about parallel edges. If three separate cables run between routers $u$
and $v$, they are three distinct edges $e_1, e_2, e_3$ with
$\partial e_i = \set{u,v}$ for each $i$. They are not ``the same edge''; a flow
algorithm can saturate one and leave the others open. Treating $E$ as a plain
subset of pairs would silently merge them, which is exactly the bug that makes a
naive max-flow implementation under-count capacity.

\begin{formalbox}
\textbf{Counting conventions.} Throughout the book, $n = \abs{V}$ and
$m = \abs{E}$. For a multigraph $m$ counts parallel edges with multiplicity. We
freely assume $V = \set{0,1,\dots,n-1}$ so vertices index into arrays in
$\Oh(1)$; relabeling an arbitrary vertex set to $\set{0,\dots,n-1}$ costs one
hash pass, $\Oh(n)$ expected, and we pay it once up front.
\end{formalbox}

\subsection{Degrees and the handshake lemma}

\begin{definition}[Degree]
\label{ch01:def-degree}
In an undirected graph the \emph{degree} $\deg(v)$ of a vertex $v$ is the number
of edges incident to it (a self-loop, when allowed, counts twice). In a digraph
the \emph{in-degree} $\deg^-(v)$ is the number of edges with head $v$ and the
\emph{out-degree} $\deg^+(v)$ is the number with tail $v$.
\end{definition}

The single most-used counting fact about graphs is also the first one anybody
proves. It is sometimes called the \emph{degree-sum formula}.

\begin{theorem}[Handshake lemma]
\label{ch01:handshake}
For any undirected multigraph $\Gph = (V,E)$,
\[
  \sum_{v \in V} \deg(v) = 2m.
\]
Consequently the number of vertices of odd degree is even.
\end{theorem}

\begin{whybox}
The idea is double counting. Build a table with one row per (vertex, incident
edge) pair and count it two ways. Summing across rows grouped by vertex gives
$\sum_v \deg(v)$. Summing grouped by edge: every edge $\set{u,v}$ touches
exactly two vertices, so it contributes exactly $2$ rows. The two totals must
agree.
\end{whybox}

\begin{proof}
Let $I = \set{(v,e) : v \in V,\ e \in E,\ v \text{ is an endpoint of } e}$ be the
set of \emph{incidences}. Count $\abs{I}$ in two ways. Fixing $v$ first, the
number of incidences with first coordinate $v$ is exactly $\deg(v)$ (a
self-loop contributes $2$, matching the degree convention), so
$\abs{I} = \sum_v \deg(v)$. Fixing $e$ first, every edge has exactly two
endpoints, so it appears in exactly two incidences, giving $\abs{I} = 2m$.
Hence $\sum_v \deg(v) = 2m$.

For the parity corollary, split the sum into even-degree and odd-degree vertices:
\[
  2m = \sum_{\deg(v) \text{ even}} \deg(v) + \sum_{\deg(v) \text{ odd}} \deg(v).
\]
The left side is even and the first right-hand sum is even, so the sum
over odd-degree vertices is even. A sum of odd numbers is even only when there
are an even count of them.
\end{proof}

The directed analogue is even simpler and just as useful: every edge has one
tail and one head, so
\[
  \sum_{v\in V}\deg^+(v) \;=\; \sum_{v\in V}\deg^-(v) \;=\; m .
\]
We will lean on these identities constantly when we bound running times, because
$\sum_v \deg(v) = 2m$ is precisely what lets ``do $\Oh(1)$ work per incident
edge at each vertex'' collapse to a clean $\Oh(n+m)$ total instead of
$\Oh(n^2)$.

\begin{labbox}
Take the five-vertex graph $V = \set{1,2,3,4,5}$ with edges
$\set{1,2},\set{1,3},\set{2,3},\set{3,4},\set{4,5}$. Read off the degrees:
$\deg(1){=}2,\ \deg(2){=}2,\ \deg(3){=}3,\ \deg(4){=}2,\ \deg(5){=}1$. Their sum
is $10 = 2\cdot 5 = 2m$. Two vertices ($3$ and $5$) have odd degree --- an even
count, as the lemma promised. This is the running example for the whole chapter;
keep it nearby.
\end{labbox}

\begin{center}
\begin{tikzpicture}[every node/.style={circle,draw=nodeblue,fill=blue!8,inner sep=2pt,minimum size=22pt,font=\small},
  edge/.style={edgegray,thick}]
  \node (1) at (0,1.4) {$1$};
  \node (2) at (0,-1.4) {$2$};
  \node (3) at (2,0) {$3$};
  \node (4) at (4,0) {$4$};
  \node (5) at (6,0) {$5$};
  \draw[edge] (1) -- (2);
  \draw[edge] (1) -- (3);
  \draw[edge] (2) -- (3);
  \draw[edge] (3) -- (4);
  \draw[edge] (4) -- (5);
\end{tikzpicture}
\end{center}

\section{Walks, paths, cycles, and connectivity}
\label{ch01:paths}

\begin{definition}[Walk, trail, path]
\label{ch01:def-walk}
A \emph{walk} of length $k$ is a sequence $v_0, e_1, v_1, e_2, \dots, e_k, v_k$
where each $e_i$ joins $v_{i-1}$ and $v_i$ (and points $v_{i-1}\to v_i$ in the
directed case). A \emph{trail} is a walk with no repeated edge. A \emph{path} is
a walk with no repeated vertex (hence no repeated edge). A walk or path is
\emph{closed} if $v_0 = v_k$; a closed path of length $\ge 1$ (in the directed
case) or length $\ge 3$ (in the simple undirected case) is a \emph{cycle}.
\end{definition}

The distinctions matter more than they look. ``There is a walk from $s$ to $t$''
and ``there is a path from $s$ to $t$'' are logically equivalent --- any walk
can be shortened to a path by deleting the loop between two visits to the same
vertex --- but algorithms produce one or the other, and proofs about lengths
care which. The next lemma gives us the basic shortening tool.

\begin{lemma}[Walk-to-path shortening]
\label{ch01:walk-to-path}
If there is a walk from $u$ to $v$, then there is a path from $u$ to $v$, and a
shortest such walk is already a path.
\end{lemma}

\begin{proof}
Take any walk $W = v_0 \dots v_k$ from $u=v_0$ to $v=v_k$. If some vertex
repeats, say $v_i = v_j$ with $i < j$, splice out the closed segment: the
shorter walk $v_0 \dots v_i v_{j+1} \dots v_k$ still goes from $u$ to $v$.
Repeat until no vertex repeats; the process terminates because each splice
strictly shortens $W$. What remains is a path. For the second claim, a
\emph{shortest} walk admits no such splice (the spliced walk would be shorter
and still valid), so it has no repeated vertex and is already a path.
\end{proof}

Geometrically, the proof excises each redundant detour loop at a repeated vertex until only a simple path remains:

\begin{center}
\begin{tikzpicture}[
    >=Stealth,
    vnode/.style={circle, draw=sourceblue, fill=sourceblue!10, thick, minimum size=7mm, font=\small\bfseries},
    repnode/.style={circle, draw=quotientorange!90!black, fill=quotientorange!15, thick, minimum size=7mm, font=\small\bfseries},
    walkedge/.style={draw=sourceblue!85!black, very thick, ->},
    loopedge/.style={draw=trapred!85!black, thick, ->, dashed},
    paneltitle/.style={font=\footnotesize\scshape, text=gray!80!black}
]
  % Top panel: Walk with a redundant cycle above v2
  \node[paneltitle, anchor=west] at (-0.5, 2.85) {Walk $W$ with self-intersection at $v_2$:};
  \node[vnode] (w0) at (0, 0.6) {$v_0$};
  \node[vnode] (w1) at (1.6, 0.6) {$v_1$};
  \node[repnode] (w2) at (3.2, 0.6) {$v_2$};
  \node[vnode] (w3) at (2.0, 2.0) {$v_3$};
  \node[vnode] (w4) at (4.4, 2.0) {$v_4$};
  \node[vnode] (w5) at (5.2, 0.6) {$v_5$};

  \draw[walkedge] (w0) -- (w1);
  \draw[walkedge] (w1) -- (w2);
  \draw[loopedge] (w2) -- (w3);
  \draw[loopedge] (w3) -- node[midway, above=3pt, font=\scriptsize, text=trapred!90!black] {redundant cycle} (w4);
  \draw[loopedge] (w4) -- (w2);
  \draw[walkedge] (w2) -- (w5);

  % Downward surgery arrow
  \draw[->, thick, gray!65] (3.2, 0.05) -- (3.2, -0.75) node[midway, right=4pt, font=\scriptsize, text=gray!75!black] {splice out cycle at $v_2$};

  % Bottom panel: Spliced simple path
  \node[paneltitle, anchor=west] at (-0.5, -1.15) {Shortened path $P$ (simple, no repeated vertices):};
  \node[vnode] (p0) at (0, -2.0) {$v_0$};
  \node[vnode] (p1) at (1.6, -2.0) {$v_1$};
  \node[repnode] (p2) at (3.2, -2.0) {$v_2$};
  \node[vnode] (p5) at (5.2, -2.0) {$v_5$};

  \draw[walkedge] (p0) -- (p1);
  \draw[walkedge] (p1) -- (p2);
  \draw[walkedge] (p2) -- (p5);
\end{tikzpicture}
\end{center}

This is why breadth-first search, which by construction explores along walks of
increasing length, returns genuine shortest \emph{paths} in unweighted graphs:
the shortest walk is automatically simple.

\begin{definition}[Connectivity]
\label{ch01:def-connectivity}
An undirected graph is \emph{connected} if every pair of vertices is joined by a
path. Its maximal connected subgraphs are its \emph{connected components}. A
digraph is \emph{strongly connected} if for every ordered pair $(u,v)$ there is
a directed path $u \to \cdots \to v$; its maximal strongly connected subgraphs
are its \emph{strongly connected components} (SCCs, the subject of
Chapter~5).
\end{definition}

``Reachable from'' --- the existence of a walk, equivalently a path --- is an
equivalence relation in the undirected case (reflexive, symmetric because edges
go both ways, transitive by concatenating walks), and its equivalence classes
are exactly the connected components. In the directed case symmetry fails, so
reachability is only a preorder, and recovering the symmetric part of it is what
makes SCCs subtler than components.

\section{Trees, forests, and DAGs}
\label{ch01:trees}

Three special shapes recur so often that they earn their own names and their own
characterization theorems.

\begin{definition}[Forest, tree]
\label{ch01:def-tree}
A \emph{forest} is an undirected graph with no cycle (it is \emph{acyclic}). A
\emph{tree} is a connected forest. A vertex of degree $1$ in a tree is a
\emph{leaf}.
\end{definition}

\begin{theorem}[Tree characterizations]
\label{ch01:tree-equiv}
For an undirected graph $T$ on $n \ge 1$ vertices, the following are equivalent.
\begin{enumerate}[label=\textup{(\roman*)}]
  \item $T$ is a tree (connected and acyclic).
  \item $T$ is connected and has exactly $n-1$ edges.
  \item $T$ is acyclic and has exactly $n-1$ edges.
  \item Every two vertices are joined by a \emph{unique} path.
\end{enumerate}
\end{theorem}

\begin{whybox}
The hub fact is: \emph{a connected graph on $n$ vertices has at least $n-1$
edges, with equality exactly when it is acyclic.} Connectivity forces enough
edges to tie $n$ vertices together; a cycle is a redundant edge you could delete
without disconnecting anything. Trees sit exactly at the boundary --- connected
but with not one edge to spare.
\end{whybox}

\begin{proof}
We show (i)$\Rightarrow$(iv)$\Rightarrow$(ii)$\Rightarrow$(iii)$\Rightarrow$(i).

(i)$\Rightarrow$(iv). Connectivity gives at least one path between any $u,v$.
Suppose there were two distinct paths $P \neq Q$ from $u$ to $v$. Let $x$ be the
last vertex of the common prefix before $P$ and $Q$ first diverge, and let $y$ be
the first vertex after $x$ along $P$ that also appears on $Q$'s suffix after $x$
(such a $y$ exists since both paths end at $v$). The segment of $P$ from $x$ to
$y$ and the segment of $Q$ from $x$ to $y$ are distinct and share no internal
vertices, so concatenating one with the reverse of the other forms a simple cycle,
contradicting acyclicity. Hence the path is unique.

(iv)$\Rightarrow$(ii). Unique paths imply connected. For the edge count, root
the tree at any vertex and map each non-root vertex to its \emph{parent edge}
(the first edge on the unique path toward the root). This is an injection from
the $n-1$ non-root vertices into $E$, and it is a surjection because every edge
$\set{x,y}$ is the parent edge of whichever of $x,y$ is farther from the root.
Hence exactly $n-1$ edges.

(ii)$\Rightarrow$(iii). Suppose $T$ is connected with $n-1$ edges but contains a
cycle. Deleting one cycle edge keeps the graph connected (the rest of the cycle
gives a detour), leaving a connected graph on $n$ vertices with $n-2$ edges. But
a connected graph needs at least $n-1$ edges: starting from $n$ singleton
components and adding edges one at a time, each edge lowers the component count
by at most one, so reaching a single component takes at least $n-1$ edges.
Contradiction, so $T$ is acyclic.

(iii)$\Rightarrow$(i). Suppose $T$ is acyclic with $n-1$ edges. Let it have $c$
components with sizes $n_1, \dots, n_c$. Each component is a connected acyclic
graph, i.e.\ a tree, so by the count just established it has $n_i - 1$ edges.
Total edges $= \sum_i (n_i - 1) = n - c$. Setting this equal to $n-1$ forces
$c = 1$: connected.
\end{proof}

\begin{corollary}[Forest edge count]
\label{ch01:forest-count}
A forest with $n$ vertices and $c$ components has exactly $n - c$ edges. Every
tree on $\ge 2$ vertices has at least two leaves.
\end{corollary}

\begin{proof}
The edge count fell out of the proof above. For leaves, if a tree had at most
one leaf then at most one vertex has degree $1$ and all others have degree
$\ge 2$ (degree $0$ is impossible in a connected graph on $\ge 2$ vertices), so
$\sum_v \deg(v) \ge 2(n-1) + 1 = 2n - 1$. By the handshake lemma this sum equals
$2(n-1) = 2n-2 < 2n-1$, a contradiction.
\end{proof}

The figure below makes the counting argument concrete on a seven-vertex tree
rooted at $r$. Each non-root vertex owns the dashed parent edge that points
toward $r$, so the six parent edges match the six non-root vertices one for one
--- the injection of part (iv)$\Rightarrow$(ii). The four leaves $\set{b,c,e,g}$
sit at the degree-$1$ fringe, comfortably more than the two the corollary
guarantees.

\begin{center}
\begin{tikzpicture}[every node/.style={circle,draw=nodeblue,fill=blue!8,inner sep=2pt,minimum size=20pt,font=\small},
  edge/.style={edgegray,thick}, pedge/.style={edgegray,thick,dashed},
  level distance=14mm, sibling distance=18mm]
  \node (r) at (3,3) {$r$};
  \node (a) at (1,1.6) {$a$};
  \node (b) at (3,1.6) {$b$};
  \node (d) at (5,1.6) {$d$};
  \node (c) at (0,0.2) {$c$};
  \node (e) at (2,0.2) {$e$};
  \node (g) at (5,0.2) {$g$};
  \draw[pedge] (a) -- (r);
  \draw[pedge] (b) -- (r);
  \draw[pedge] (d) -- (r);
  \draw[pedge] (c) -- (a);
  \draw[pedge] (e) -- (a);
  \draw[pedge] (g) -- (d);
\end{tikzpicture}
\end{center}

Now the directed cousin.

\begin{definition}[DAG]
\label{ch01:def-dag}
A \emph{directed acyclic graph} (DAG) is a digraph with no directed cycle. A
\emph{topological order} of a digraph is a linear ordering of $V$ such that every
edge $(u,v)$ has $u$ before $v$.
\end{definition}

\begin{theorem}[DAG iff topologically orderable]
\label{ch01:dag-topo}
A digraph admits a topological order if and only if it is a DAG.
\end{theorem}

\begin{proof}
($\Rightarrow$) If a topological order exists, any directed cycle
$v_0 \to v_1 \to \cdots \to v_k = v_0$ would force $v_0$ before $v_1$ before
$\cdots$ before $v_0$ in the order, impossible. So no cycle: a DAG.

($\Leftarrow$) Suppose $\Gph$ is acyclic. We claim some vertex has in-degree
$0$. If not, every vertex has an incoming edge; follow edges backward from any
start, and after $n$ steps you have visited $n+1$ vertices, so two coincide,
producing a directed cycle --- contradiction. Place a source $s$ (in-degree $0$)
first, delete it, and recurse on the remaining graph; deleting a vertex cannot
create a cycle, so the recursion stays valid. This builds a full topological
order.
\end{proof}

The constructive direction is exactly Kahn's algorithm, which Chapter~4
turns into linear-time pseudocode. The picture below shows the same five
vertices oriented into a DAG and one valid topological order written left to
right.

\begin{center}
\begin{tikzpicture}[every node/.style={circle,draw=nodeblue,fill=blue!8,inner sep=2pt,minimum size=22pt,font=\small},
  >={Stealth[length=5pt]}, edge/.style={edgegray,thick,->}]
  \node (1) at (0,0) {$1$};
  \node (2) at (1.8,0.9) {$2$};
  \node (3) at (3.6,0) {$3$};
  \node (4) at (5.4,0) {$4$};
  \node (5) at (7.2,0) {$5$};
  \draw[edge] (1) -- (2);
  \draw[edge] (1) -- (3);
  \draw[edge] (2) -- (3);
  \draw[edge] (3) -- (4);
  \draw[edge] (4) -- (5);
  \node[draw=none,fill=none,font=\small] at (3.6,-1.2) {topological order: $1,\,2,\,3,\,4,\,5$};
\end{tikzpicture}
\end{center}

\section{Bipartite graphs}
\label{ch01:bipartite}

\begin{definition}[Bipartite]
\label{ch01:def-bipartite}
A graph is \emph{bipartite} if its vertices can be split into two sets $X$ and
$Y$ (the \emph{sides}) so that every edge has one endpoint in $X$ and the other
in $Y$. Equivalently, the vertices can be $2$-colored so no edge is
monochromatic.
\end{definition}

Bipartite graphs model anything with a natural two-sided structure: jobs and
workers, students and courses, the two halves of a matching problem. The clean
test for bipartiteness is one of the prettiest small theorems in the book.

\begin{theorem}[Odd-cycle characterization]
\label{ch01:bip-oddcycle}
A graph is bipartite if and only if it contains no cycle of odd length.
\end{theorem}

\begin{whybox}
Walk around a cycle alternating colors $X, Y, X, Y, \dots$. To come back to your
starting color you must take an even number of steps. So an odd cycle can never
be properly $2$-colored, and conversely if there is no odd cycle the alternating
coloring never contradicts itself.
\end{whybox}

\begin{proof}
($\Rightarrow$) In a $2$-coloring the colors alternate along any walk, so a
closed walk returns to its start color only after evenly many steps. Every cycle
is a closed walk, hence even.

($\Leftarrow$) Assume no odd cycle; we $2$-color each connected component
independently. Fix a component, pick a root $r$, and color each vertex $v$ by
the parity of $\dist(r,v)$ (the length of a shortest path, well defined within a
connected component). Suppose for contradiction some edge $\set{u,v}$ is
monochromatic, i.e.\ $\dist(r,u)$ and $\dist(r,v)$ have the same parity. Take
shortest paths $P_u$ from $r$ to $u$ and $P_v$ from $r$ to $v$; let $z$ be the
last vertex on $P_u$ that also lies on $P_v$. Any subpath of a shortest path is
itself shortest, so the prefixes to $z$ on both paths have length $\dist(r,z)$,
and the suffixes $z\rightsquigarrow u$ along $P_u$ and $z\rightsquigarrow v$
along $P_v$ share no vertices other than $z$. Those two suffixes have lengths
$\dist(r,u)-\dist(r,z)$ and $\dist(r,v)-\dist(r,z)$, which have the same parity,
so their sum is even; adding the edge $\set{u,v}$ closes a simple cycle of
\emph{odd} length. Contradiction. So no edge is monochromatic and the graph is
bipartite.
\end{proof}

This proof is constructive: run a BFS, color by distance parity, and check every
edge. That is a linear-time bipartiteness test, and we will reuse the exact idea
when BFS gets its own chapter. The two panels below contrast a valid $2$-coloring
with the odd cycle closed whenever an edge $\set{u,v}$ joins two vertices of the
same distance parity from $r$:

\begin{center}
\begin{tikzpicture}[
    >=Stealth,
    vnode/.style={circle, draw=nodeblue, thick, inner sep=2pt, minimum size=20pt, font=\small},
    edge/.style={edgegray, thick},
    xv/.style={vnode, fill=sourceblue!22},
    yv/.style={vnode, fill=targetgreen!22},
    paneltitle/.style={font=\footnotesize\scshape, text=gray!80!black}
]
  % Left panel: Bipartite 2-coloring
  \node[paneltitle] at (1.5, 2.8) {Bipartite ($2$-colorable)};
  \node[xv] (x1) at (0, 2.0) {$x_1$};
  \node[xv] (x2) at (0, 1.0) {$x_2$};
  \node[xv] (x3) at (0, 0.0) {$x_3$};
  \node[yv] (y1) at (3.0, 1.5) {$y_1$};
  \node[yv] (y2) at (3.0, 0.5) {$y_2$};
  \draw[edge] (x1) -- (y1);
  \draw[edge] (x2) -- (y1);
  \draw[edge] (x2) -- (y2);
  \draw[edge] (x3) -- (y2);
  \node[font=\footnotesize, text=sourceblue!90!black] at (0, -0.75) {side $X$ (even)};
  \node[font=\footnotesize, text=targetgreen!90!black] at (3.0, -0.75) {side $Y$ (odd)};

  % Right panel: Odd-cycle witness from monochromatic edge {u,v}
  \node[paneltitle] at (8.2, 2.8) {Same-parity edge $\set{u,v}$ closes odd cycle};
  \node[xv] (r) at (6.0, 1.0) {$r$};
  \node[yv] (z) at (7.6, 1.0) {$z$};
  \node[xv] (u) at (9.6, 1.9) {$u$};
  \node[xv] (v) at (9.6, 0.1) {$v$};
  \draw[edge, dashed] (r) -- node[midway, above=2pt, font=\scriptsize] {$P_z$} (z);
  \draw[edge, sourceblue!85!black, very thick] (z) -- node[midway, above left=1pt, font=\scriptsize] {even/odd} (u);
  \draw[edge, sourceblue!85!black, very thick] (z) -- node[midway, below left=1pt, font=\scriptsize] {same parity} (v);
  \draw[trapred!85!black, very thick] (u) -- node[midway, right=3pt, font=\scriptsize, text=trapred!90!black] {conflict $\set{u,v}$} (v);
  \node[font=\footnotesize, text=trapred!90!black] at (8.4, -0.75) {two equal-parity branches $+ 1$ edge $=$ odd cycle};
\end{tikzpicture}
\end{center}

\section{Three ways to store a graph}
\label{ch01:representations}

We now leave the mathematics and confront the engineering. The same abstract
graph $(V,E)$ can be packed into memory in several ways, and the choice governs
how fast every later algorithm runs. The three standard structures are the
\emph{adjacency matrix}, the \emph{adjacency list}, and the \emph{edge list}.
We measure each on four operations that algorithms perform constantly:

\begin{itemize}
  \item \textbf{space}: total memory;
  \item \textbf{edge query} \texttt{has\_edge(u,v)}: is $(u,v)$ present?
  \item \textbf{neighbor scan}: enumerate all neighbors of a given $v$;
  \item \textbf{global scan}: enumerate all edges once.
\end{itemize}

We run all three on the running five-vertex graph
$\set{1,2},\set{1,3},\set{2,3},\set{3,4},\set{4,5}$ so you can see the bytes.

\subsection{Adjacency matrix}

\begin{definition}[Adjacency matrix]
\label{ch01:def-adjmatrix}
The \emph{adjacency matrix} of $\Gph$ on vertex set $\set{0,\dots,n-1}$ is the
$n\times n$ matrix $\Adjm$ with $\Adjm_{uv} = 1$ if $(u,v) \in E$ and $0$
otherwise (for a weighted graph, store $w(u,v)$ in place of $1$, with a sentinel
such as $+\infty$ or $0$ for absence). For an undirected graph $\Adjm$ is
symmetric; for a multigraph store the multiplicity.
\end{definition}

For our graph (with vertices $1\dots5$ indexing rows and columns), the
undirected adjacency matrix is

\[
\Adjm =
\begin{pmatrix}
0 & 1 & 1 & 0 & 0 \\
1 & 0 & 1 & 0 & 0 \\
1 & 1 & 0 & 1 & 0 \\
0 & 0 & 1 & 0 & 1 \\
0 & 0 & 0 & 1 & 0
\end{pmatrix}.
\]

Notice the symmetry, and notice the row sums: row $i$ sums to $\deg(i)$, and the
grand total of all entries is $\sum_i \deg(i) = 2m = 10$, the handshake lemma
showing up as a matrix identity.

\begin{focusbox}
The adjacency matrix turns graph questions into linear algebra. The entry
$(\Adjm^k)_{uv}$ counts walks of length exactly $k$ from $u$ to $v$ (prove it by
induction: $(\Adjm^{k})_{uv} = \sum_z (\Adjm^{k-1})_{uz}\Adjm_{zv}$ is exactly
``length-$(k-1)$ walks to a neighbor of $v$''). This observation is the
starting point for the spectral half of the book --- eigenvalues of $\Adjm$ and of
the Laplacian $\Lap = \mathbf{D} - \Adjm$ control mixing, expansion, and
clustering. See Chapters~31--35.
\end{focusbox}

\textbf{Costs.} Space is $\Th(n^2)$: every one of the $n^2$ cells is stored
whether or not the edge exists. The edge query is where the matrix excels:
\texttt{has\_edge(u,v)} is a single array lookup, $\Oh(1)$. A neighbor
scan, though, is $\Th(n)$ regardless of how few neighbors $v$ has --- you must
walk the entire row testing each cell. A global edge scan is $\Th(n^2)$.

\subsection{Adjacency list}

\begin{definition}[Adjacency list]
\label{ch01:def-adjlist}
The \emph{adjacency list} stores, for each vertex $v$, a list
$\mathrm{Adj}[v]$ of its out-neighbors (in the undirected case, of its
neighbors; each undirected edge $\set{u,v}$ appears twice, once in
$\mathrm{Adj}[u]$ and once in $\mathrm{Adj}[v]$). For a weighted graph each list
entry carries the weight alongside the neighbor.
\end{definition}

For our graph:
\[
\begin{aligned}
\mathrm{Adj}[1] &= [\,2,\,3\,] & \mathrm{Adj}[2] &= [\,1,\,3\,] &
\mathrm{Adj}[3] &= [\,1,\,2,\,4\,] \\
\mathrm{Adj}[4] &= [\,3,\,5\,] & \mathrm{Adj}[5] &= [\,4\,] .
\end{aligned}
\]

The total number of list entries is $\sum_v \deg(v) = 2m = 10$: the handshake
lemma is now literally the size of the data structure. That is the main reason
adjacency lists are the default representation in this book. Drawn as an array of
per-vertex headers pointing to their neighbor lists, the five vertices and ten
directed incidences look like this:

\begin{center}
\begin{tikzpicture}[
  vbox/.style={draw=nodeblue,fill=blue!8,minimum size=20pt,font=\small},
  cell/.style={draw=edgegray,fill=white,minimum width=18pt,minimum height=18pt,font=\small},
  >={Stealth[length=4pt]}]
  \foreach \v/\y in {1/4,2/3,3/2,4/1,5/0} {
    \node[vbox] (v\v) at (0,\y) {$\v$};
  }
  \node[cell,right=6mm of v1] (a11) {2}; \node[cell,right=0 of a11] (a12) {3};
  \node[cell,right=6mm of v2] (a21) {1}; \node[cell,right=0 of a21] (a22) {3};
  \node[cell,right=6mm of v3] (a31) {1}; \node[cell,right=0 of a31] (a32) {2}; \node[cell,right=0 of a32] (a33) {4};
  \node[cell,right=6mm of v4] (a41) {3}; \node[cell,right=0 of a41] (a42) {5};
  \node[cell,right=6mm of v5] (a51) {4};
  \draw[->] (v1) -- (a11); \draw[->] (v2) -- (a21); \draw[->] (v3) -- (a31);
  \draw[->] (v4) -- (a41); \draw[->] (v5) -- (a51);
\end{tikzpicture}
\end{center}

\textbf{Costs.} Space is $\Th(n + m)$: one array slot per vertex plus one list
entry per edge-endpoint. A neighbor scan of $v$ is $\Th(1 + \deg(v))$ --- you
touch exactly the neighbors and nothing else, which is optimal. A global edge
scan visits every list once for a total of $\Th(n + m)$. The weak spot is the
edge query: \texttt{has\_edge(u,v)} must search $\mathrm{Adj}[u]$, costing
$\Oh(\deg(u))$ in the worst case (or $\Oh(\log\deg(u))$ if the list is kept
sorted, or expected $\Oh(1)$ if you maintain a parallel hash set of edges).

\begin{trapbox}
The two most common adjacency-list bugs are both about \emph{the two copies of
an undirected edge}. (1) Forgetting to add the second copy: you push $v$ into
$\mathrm{Adj}[u]$ but not $u$ into $\mathrm{Adj}[v]$, so BFS from $v$ silently
cannot reach $u$. (2) For algorithms that walk both copies --- Euler tours,
many flow networks --- you must find an edge's \emph{twin} in $\Oh(1)$. The
standard trick (used throughout the flow chapters) is to store edges in one
array and put copies at indices $2i$ and $2i{+}1$, so the twin of edge $e$ is
$e \oplus 1$. Get this wrong and your residual graph leaks capacity.
\end{trapbox}

\subsection{Edge list}

\begin{definition}[Edge list]
\label{ch01:def-edgelist}
The \emph{edge list} is simply the array of edges, each stored as a triple
$(u, v, w)$ (or pair $(u,v)$ when unweighted). No per-vertex indexing.
\end{definition}

For our graph: $[(1,2),(1,3),(2,3),(3,4),(4,5)]$. That is the whole structure.

\textbf{Costs.} Space $\Th(n + m)$ (the $n$ is only to know the vertex count).
There is no fast neighbor scan or edge query without preprocessing --- both are
$\Th(m)$ linear scans. What the edge list is unbeatable at is \emph{processing
every edge once, in an order you control}. Kruskal's MST algorithm sorts the
edge list by weight and sweeps it; Bellman--Ford relaxes the entire edge list
$n-1$ times. For those algorithms the edge list is the most direct representation.

\begin{historybox}
The trade-off between matrix and list is older than computer science. The
incidence and adjacency matrices go back to the birth of algebraic graph theory
in the nineteenth century --- Kirchhoff's 1847 work on electrical networks
already manipulates what we now call the incidence matrix, and we will meet it
again in the Laplacian chapters. Adjacency lists as a deliberate
\emph{algorithmic} choice --- selected precisely so that running time scales with
$m$ rather than $n^2$ --- crystallized in the 1970s with Hopcroft and Tarjan,
whose linear-time DFS-based algorithms for planarity and connectivity simply do
not work in $\Oh(n^2)$ and forced the field to take representation seriously.
\end{historybox}

\subsection{The comparison, side by side}

\begin{center}
\renewcommand{\arraystretch}{1.3}
\begin{tabular}{lccc}
\toprule
 & \textbf{Adjacency matrix} & \textbf{Adjacency list} & \textbf{Edge list} \\
\midrule
space            & $\Th(n^2)$ & $\Th(n+m)$ & $\Th(n+m)$ \\
\texttt{has\_edge(u,v)} & $\Oh(1)$ & $\Oh(\deg u)$ & $\Oh(m)$ \\
scan nbrs of $v$ & $\Th(n)$ & $\Th(1+\deg v)$ & $\Th(m)$ \\
scan all edges   & $\Th(n^2)$ & $\Th(n+m)$ & $\Th(m)$ \\
add an edge      & $\Oh(1)$ & $\Oh(1)$ (push) & $\Oh(1)$ (push) \\
delete an edge   & $\Oh(1)$ & $\Oh(\deg)$ & $\Oh(m)$ \\
\bottomrule
\end{tabular}
\end{center}

\begin{intuitionbox}
A graph is \emph{dense} when $m = \Th(n^2)$ and \emph{sparse} when
$m = \Oh(n)$ or so. On a dense graph the matrix's $\Th(n^2)$ space is no worse
than the list's $\Th(n+m) = \Th(n^2)$, and the $\Oh(1)$ edge query is a real
win, so matrices shine. On a sparse graph --- which is almost every graph that
arises in practice, from road maps to social networks to web crawls --- the
matrix wastes nearly all of its $n^2$ cells on zeros, and the adjacency list's
$\Th(n+m)$ is dramatically smaller. \textbf{Default to adjacency lists.} Reach
for a matrix only when the graph is dense or when an algorithm genuinely needs
$\Oh(1)$ random-access edge tests (Floyd--Warshall, some matrix-multiplication
shortest-path tricks).
\end{intuitionbox}

\subsection{Building an adjacency list in linear time}

When graph algorithms meet real hardware, nested pointer-chasing adjacency lists
(such as nested vectors or pointer-linked chains) incur significant memory
overhead and cache misses. Instead, high-performance graph processing packs the entire
graph into two flat, contiguous arrays: the \emph{Compressed Sparse Row} (CSR)
representation. An arbitrary edge list can be converted into CSR layout
in $\Th(n+m)$ time without sorting or hashing, using a two-pass counting prefix sum.

\begin{algorithm}[H]
\caption{\textsc{BuildCSR}: edge list to compressed adjacency list}
\label{ch01:alg-csr}
\KwIn{$n$, edge list $E = [(u_1,v_1),\dots,(u_m,v_m)]$ (directed; for undirected, add both $(u,v)$ and $(v,u)$).}
\KwOut{arrays $\mathrm{head}[0..n]$ and $\mathrm{dest}[0..m-1]$ with $\mathrm{Adj}[v] = \mathrm{dest}[\mathrm{head}[v]\,..\,\mathrm{head}[v{+}1]{-}1]$.}
\For{$v \gets 0$ \KwTo $n$}{$\mathrm{head}[v] \gets 0$\;}
\For{$i \gets 1$ \KwTo $m$}{$\mathrm{head}[u_i + 1] \gets \mathrm{head}[u_i + 1] + 1$\;} \tcp{count out-degrees}
\For{$v \gets 1$ \KwTo $n$}{$\mathrm{head}[v] \gets \mathrm{head}[v] + \mathrm{head}[v-1]$\;} \tcp{prefix sums = block starts}
$\mathrm{cur} \gets \mathrm{copy}(\mathrm{head})$\;
\For{$i \gets 1$ \KwTo $m$}{
  $\mathrm{dest}[\mathrm{cur}[u_i]] \gets v_i$\;
  $\mathrm{cur}[u_i] \gets \mathrm{cur}[u_i] + 1$\;
}
\Return $(\mathrm{head}, \mathrm{dest})$\;
\end{algorithm}

\begin{whybox}
\textbf{Why it is correct.} After the counting loop, $\mathrm{head}[u+1]$ holds
$\deg^+(u)$. The prefix-sum loop turns these counts into starting offsets:
$\mathrm{head}[v]$ becomes the index where $v$'s block begins, because the blocks
are laid out in vertex order and each occupies $\deg^+$ slots. The placement
loop uses $\mathrm{cur}[u]$ as a moving write-cursor into $u$'s block; since we
increment it after each write and the block has exactly $\deg^+(u)$ slots, the
$\deg^+(u)$ edges out of $u$ fill the block exactly, never overflowing into the
next vertex's block. Distinct vertices write to disjoint blocks, so nothing
collides. \textbf{Why it is fast.} Four loops, each $\Th(n)$ or $\Th(m)$, with
$\Oh(1)$ work per iteration: $\Th(n+m)$ total, no comparisons, no hashing.
\end{whybox}

\begin{labbox}
Trace \textsc{BuildCSR} on the directed edge list
$[(0,1),(0,2),(1,2),(2,3),(3,1)]$ with $n=4$. Counting gives out-degrees
$\deg^+ = [2,1,1,1]$, so after the count loop
$\mathrm{head} = [0,2,1,1,1]$ (slot $v{+}1$ holds $\deg^+(v)$). Prefix-summing:
$\mathrm{head} = [0,2,3,4,5]$. Now place: edge $(0,1)$ writes $\mathrm{dest}[0]{=}1$,
cursor of $0$ moves to $1$; $(0,2)$ writes $\mathrm{dest}[1]{=}2$; $(1,2)$ writes
$\mathrm{dest}[2]{=}2$; $(2,3)$ writes $\mathrm{dest}[3]{=}3$; $(3,1)$ writes
$\mathrm{dest}[4]{=}1$. Final $\mathrm{dest} = [1,2,2,3,1]$, and indeed
$\mathrm{Adj}[0]=\mathrm{dest}[0..1]=[1,2]$ and $\mathrm{Adj}[2]=\mathrm{dest}[3..3]=[3]$.
Check it against the edges.
\end{labbox}

The resulting CSR layout packs all five directed edges contiguously into
$\mathrm{dest}$, with each $\mathrm{head}[v]$ indexing the start of vertex $v$'s
slice:

\begin{center}
\begin{tikzpicture}[
    >=Stealth,
    vnode/.style={circle, draw=sourceblue, fill=sourceblue!10, thick, minimum size=6.5mm, font=\small\bfseries},
    dedge/.style={draw=edgegray, thick, ->},
    cell/.style={draw=gray!70, thick, rounded corners=2pt, minimum width=9.5mm, minimum height=7.5mm, font=\small\ttfamily},
    ruler/.style={font=\tiny\ttfamily, text=gray!80!black},
    paneltitle/.style={font=\footnotesize\scshape, text=gray!80!black},
    labelcol/.style={font=\footnotesize\bfseries\ttfamily, text=gray!90!black, anchor=east}
]
  % Top Lane: The concrete 4-vertex directed graph
  \node[paneltitle, anchor=east] at (1.5, 2.8) {Graph $G$ ($n{=}4, m{=}5$):};
  \node[vnode, fill=sourceblue!20] (g0) at (2.5, 2.8) {$0$};
  \node[vnode, fill=targetgreen!20] (g1) at (5.0, 3.5) {$1$};
  \node[vnode, fill=quotientorange!20] (g2) at (5.0, 2.1) {$2$};
  \node[vnode, fill=structurepurple!20] (g3) at (7.5, 2.8) {$3$};

  \draw[dedge] (g0) -- (g1);
  \draw[dedge] (g0) -- (g2);
  \draw[dedge] (g1) -- (g2);
  \draw[dedge] (g2) -- (g3);
  \draw[dedge] (g3) to[bend right=25] (g1);

  % Middle Lane: head array (n+1 = 5 slots)
  \node[labelcol] at (1.3, 0.7) {head:};
  \foreach \i/\x in {0/2.0, 1/3.6, 2/5.2, 3/6.8, 4/8.4} {
    \node[ruler] at (\x, 1.3) {\i};
  }
  \node[cell, fill=sourceblue!15, draw=sourceblue!70] (h0) at (2.0, 0.7) {0};
  \node[cell, fill=targetgreen!15, draw=targetgreen!70] (h1) at (3.6, 0.7) {2};
  \node[cell, fill=quotientorange!15, draw=quotientorange!70] (h2) at (5.2, 0.7) {3};
  \node[cell, fill=structurepurple!15, draw=structurepurple!70] (h3) at (6.8, 0.7) {4};
  \node[cell, fill=gray!12, draw=gray!60] (h4) at (8.4, 0.7) {5};
  \node[font=\tiny, text=gray!70!black, right=2pt of h4] {(sentinel)};

  % Bottom Lane: dest array (m = 5 slots)
  \node[labelcol] at (1.3, -1.0) {dest:};
  \node[cell, fill=sourceblue!15, draw=sourceblue!70] (d0) at (2.0, -1.0) {1};
  \node[cell, fill=sourceblue!15, draw=sourceblue!70] (d1) at (3.6, -1.0) {2};
  \node[cell, fill=targetgreen!15, draw=targetgreen!70] (d2) at (5.2, -1.0) {2};
  \node[cell, fill=quotientorange!15, draw=quotientorange!70] (d3) at (6.8, -1.0) {3};
  \node[cell, fill=structurepurple!15, draw=structurepurple!70] (d4) at (8.4, -1.0) {1};
  \foreach \i/\x in {0/2.0, 1/3.6, 2/5.2, 3/6.8, 4/8.4} {
    \node[ruler] at (\x + 0.34, -0.45) {\i};
  }

  % Connecting pointer arrows from head to dest
  \draw[->, thick, sourceblue!85!black] (h0.south) -- (d0.north);
  \draw[->, thick, targetgreen!85!black] (h1.south) -- (d2.north);
  \draw[->, thick, quotientorange!85!black] (h2.south) -- (d3.north);
  \draw[->, thick, structurepurple!85!black] (h3.south) -- (d4.north);

  % Brackets or slices beneath dest
  \node[font=\scriptsize, text=sourceblue!90!black] at (2.8, -1.75) {$\underbrace{\hspace{2.3cm}}_{\mathrm{Adj}[0]=\{1,2\}}$};
  \node[font=\scriptsize, text=targetgreen!90!black] at (5.2, -1.75) {$\underbrace{\hspace{0.9cm}}_{\mathrm{Adj}[1]=\{2\}}$};
  \node[font=\scriptsize, text=quotientorange!90!black] at (6.8, -1.75) {$\underbrace{\hspace{0.9cm}}_{\mathrm{Adj}[2]=\{3\}}$};
  \node[font=\scriptsize, text=structurepurple!90!black] at (8.4, -1.75) {$\underbrace{\hspace{0.9cm}}_{\mathrm{Adj}[3]=\{1\}}$};
\end{tikzpicture}
\end{center}

\section{Implicit graphs}
\label{ch01:implicit}

Sometimes the graph is never stored at all. An \emph{implicit graph} is given by
a starting vertex and a function that, asked about any vertex $v$, returns its
neighbors --- a successor oracle $\mathrm{succ}(v)$ --- without any global edge
list ever materializing.

\begin{intuitionbox}
Think of a Rubik's cube. The ``vertices'' are the $4.3\times 10^{19}$
configurations; the ``edges'' connect configurations one twist apart. Nobody
stores that graph --- it would not fit on every hard drive on Earth. But you can
still search it, because from any configuration you can \emph{generate} the dozen
or so neighbors on demand. The graph exists only as a rule, and BFS or A*
explores it one expanded vertex at a time.
\end{intuitionbox}

\begin{definition}[Implicit graph]
\label{ch01:def-implicit}
An \emph{implicit graph} is specified by a vertex type, a designated start
vertex (or set), and a successor function $\mathrm{succ}: V \to 2^V$ computable
on demand. The graph $(V,E)$ it represents has
$E = \set{(v,u) : u \in \mathrm{succ}(v)}$, but neither $V$ nor $E$ is held in
memory; the search visits only the reachable portion it actually touches.
\end{definition}

Grid pathfinding is the everyday example: the vertices are cells $(r,c)$ of a
maze, and $\mathrm{succ}((r,c))$ returns the up-to-four orthogonal neighbors that
are in-bounds and not walls. You would never build an adjacency list; you
compute neighbors arithmetically.

\begin{algorithm}[H]
\caption{\textsc{ImplicitBFS}: shortest path on an on-demand graph}
\label{ch01:alg-implicit-bfs}
\KwIn{start $s$, goal-test $\mathrm{isGoal}$, successor oracle $\mathrm{succ}$.}
\KwOut{the length of a shortest path from $s$ to a goal, or $\infty$.}
$Q \gets$ empty queue\; $\mathrm{seen} \gets \set{s}$\; enqueue $(s, 0)$\;
\While{$Q$ nonempty}{
  $(v, d) \gets$ dequeue\;
  \If{$\mathrm{isGoal}(v)$}{\Return $d$\;}
  \ForEach{$u \in \mathrm{succ}(v)$}{
    \If{$u \notin \mathrm{seen}$}{
      add $u$ to $\mathrm{seen}$\; enqueue $(u, d+1)$\;
    }
  }
}
\Return $\infty$\;
\end{algorithm}

Exploring outward level by level --- the principle behind Breadth-First Search,
which we develop rigorously in Chapter~2 --- works directly on implicit graphs.
The only structural adaptation is that $\mathrm{seen}$ is maintained as a hash set
rather than a length-$n$ array, since $n$ is unknown and the reachable portion
may be a tiny fraction of the state space. Just as in standard BFS, states leave
the queue in nondecreasing order of distance from $s$, so the first time a goal
is dequeued its associated distance $d$ is the true shortest-path length.
The runtime is $\Oh(\abs{V'} + \abs{E'})$, where $V'$ and $E'$ are the vertices
and transitions actually explored --- often exponentially smaller than the full
state space $(V, E)$. The figure below shows this on a $3\times 3$ grid with two
blocked cells: the successor oracle $\mathrm{succ}((r,c))$ generates only valid
open neighbors on demand, materializing a small search tree without ever
allocating a graph structure for the unvisited region:

\begin{center}
\begin{tikzpicture}[
    >=Stealth,
    gcell/.style={draw=gray!70, thick, minimum width=9.5mm, minimum height=9.5mm, font=\scriptsize\ttfamily},
    wall/.style={gcell, fill=gray!45, text=white},
    vnode/.style={circle, draw=sourceblue, fill=sourceblue!12, thick, minimum size=6.8mm, font=\scriptsize\ttfamily},
    paneltitle/.style={font=\footnotesize\scshape, text=gray!80!black}
]
  % Left panel: 3x3 grid maze
  \node[paneltitle] at (0.95, 2.75) {Implicit grid rule $\mathrm{succ}(r,c)$};
  \node[gcell, fill=sourceblue!22, draw=sourceblue!80!black] at (0, 1.9) {s:0};
  \node[gcell, fill=sourceblue!10] at (0.95, 1.9) {d=1};
  \node[gcell, fill=sourceblue!10] at (1.9, 1.9) {d=2};

  \node[gcell, fill=sourceblue!10] at (0, 0.95) {d=1};
  \node[wall] at (0.95, 0.95) {wall};
  \node[gcell, fill=targetgreen!22, draw=targetgreen!80!black] at (1.9, 0.95) {g:3};

  \node[gcell, fill=sourceblue!10] at (0, 0) {d=2};
  \node[wall] at (0.95, 0) {wall};
  \node[gcell, fill=gray!6, text=gray!70] at (1.9, 0) {unseen};

  % Right panel: Materialized BFS tree V', E'
  \node[paneltitle] at (7.2, 2.75) {Materialized subgraph $(V',E')$ on demand};
  \node[vnode, fill=sourceblue!22] (s00) at (4.6, 0.95) {0,0};
  \node[vnode] (s01) at (6.4, 1.75) {0,1};
  \node[vnode] (s10) at (6.4, 0.15) {1,0};
  \node[vnode] (s02) at (8.2, 1.75) {0,2};
  \node[vnode] (s20) at (8.2, 0.15) {2,0};
  \node[vnode, fill=targetgreen!22, draw=targetgreen!80!black] (s12) at (10.0, 1.75) {1,2};

  \draw[->, thick, sourceblue!85!black] (s00) -- (s01);
  \draw[->, thick, sourceblue!85!black] (s00) -- (s10);
  \draw[->, thick, sourceblue!85!black] (s01) -- (s02);
  \draw[->, thick, sourceblue!85!black] (s10) -- (s20);
  \draw[->, thick, targetgreen!80!black, very thick] (s02) -- node[midway, above=2pt, font=\scriptsize, text=targetgreen!85!black] {goal!} (s12);
\end{tikzpicture}
\end{center}

\begin{applicationbox}
Implicit graphs are how search powers puzzle solvers (sliding tiles, Rubik's
cube), AI planning (states and actions), game trees (positions and moves), and
model checking (program states and transitions). The same BFS/DFS/Dijkstra
machinery from the next chapters applies verbatim; you only swap the explicit
adjacency structure for a $\mathrm{succ}$ oracle. The connection runs the other
way too: every explicit graph \emph{is} an implicit graph whose oracle is ``look
up $\mathrm{Adj}[v]$.''
\end{applicationbox}

\section{The cost model}
\label{ch01:costmodel}

Every running time in this book is a function of two numbers:
\[
  n = \abs{V} \quad\text{(vertices)}, \qquad m = \abs{E} \quad\text{(edges)}.
\]
We state complexities in terms of both because the relationship between them is
where all the interesting trade-offs live. Three facts anchor the whole scale.

\begin{proposition}[Edge-count bounds]
\label{ch01:edge-bounds}
For a simple undirected graph on $n$ vertices, $0 \le m \le \binom{n}{2} =
\frac{n(n-1)}{2}$. For a simple digraph, $0 \le m \le n(n-1)$. If the graph is
connected, $m \ge n-1$.
\end{proposition}

\begin{proof}
The upper bounds count the available vertex pairs (unordered for undirected,
ordered for directed; no self-loops). The lower bound for connected graphs is
the spanning-tree fact from Theorem~\ref{ch01:tree-equiv}: a connected graph
contains a spanning tree, which already has $n-1$ edges.
\end{proof}

\begin{focusbox}
\textbf{Why we write $\Oh(n+m)$ and not $\Oh(m)$.} On a connected graph
$m \ge n-1$, so $n + m = \Th(m)$ and the two are interchangeable. But a graph
can have isolated vertices and many components, where $m$ can be as small as $0$
while $n$ is large. An algorithm that must at least look at every vertex (BFS
initializing a visited array, say) genuinely costs $\Th(n+m)$, and writing
$\Oh(m)$ would be wrong when $m \ll n$. We keep both terms so the bound is
correct on \emph{every} input, dense or sparse, connected or shattered.
\end{focusbox}

\begin{trapbox}
``Linear time'' for a graph algorithm means $\Oh(n+m)$, not $\Oh(n)$ and not
$\Oh(\text{input length})$ measured in some other unit. An algorithm touching
every edge once is linear; one touching every \emph{pair} of vertices is
$\Th(n^2)$ and is linear only on dense graphs. Quote $n$ and $m$ separately and
the distinction never bites you.
\end{trapbox}

\begin{complexitybox}
\textbf{The yardsticks used throughout the book.}
\begin{itemize}
  \item $\Oh(n+m)$ --- ``graph-linear.'' BFS, DFS, topological sort, SCCs,
        bridges. The gold standard: you cannot beat reading the graph once.
  \item $\Oh(m\log n)$ --- Dijkstra with a binary min-heap, Kruskal/Prim. The
        factor $\log n$ reflects the tree height $\lfloor \log_2 n \rfloor$ of a
        binary heap: extracting the minimum-distance vertex or decreasing a key
        traverses at most one root-to-leaf branch.
  \item $\Oh(m + n\log n)$ --- Dijkstra with a Fibonacci heap. By maintaining a
        flexible collection of trees with lazy consolidation, Fibonacci heaps achieve
        amortized $\Oh(1)$ \texttt{decrease-key} operations, so the $m$ edge
        relaxations run in linear time while only the $n$ minimum-extractions pay
        the $\Oh(\log n)$ consolidation cost.
  \item $\Oh(nm)$ --- Bellman--Ford, Edmonds--Karp max-flow. Polynomial but the
        step up from linear is steep.
  \item $\Oh(n^3)$ --- Floyd--Warshall all-pairs, naive matrix methods. Fine for
        dense graphs of moderate $n$.
  \item $\widetilde{\Oh}(m)$ --- the modern frontier: near-linear Laplacian
        solvers and electrical-flow methods (Chapter~41), where $\widetilde{\Oh}$
        hides polylogarithmic factors.
\end{itemize}
We assume the word-RAM model: vertex labels, edge weights, and array indices fit
in a machine word, and arithmetic and array access cost $\Oh(1)$.
\end{complexitybox}

\begin{connectionbox}
The choice of representation matters throughout the rest of the book and in its
companion volume on linear algebra. The adjacency matrix $\Adjm$ and its
relative the Laplacian $\Lap = \mathbf{D} - \Adjm$ connect graphs to linear algebra:
a graph question becomes a matrix question, and the spectrum of $\Lap$ encodes
connectivity, expansion, and random walks (Chapters~31--36). The sparse
adjacency list, meanwhile, makes near-linear algorithms conceivable: operations
scan only existing neighbors rather than empty entries.

More immediately, Chapter~2 uses the adjacency list in our first algorithm,
Breadth-First Search. We explore the graph level by level, prove the
queue invariant that guarantees exact shortest paths on unweighted graphs, and
verify that adjacency lists deliver the optimal $\Oh(n+m)$ running time.
\end{connectionbox}

\section{Exercises}
\label{ch01:exercises}

\begin{exercise}[Handshake in practice]
\label{ch01:ex-handshake}
A graph has degree sequence $(5,3,3,2,2,1)$. Is such a simple graph possible?
Either draw one or prove it cannot exist. (Hint: first check the handshake
parity, then think about whether the vertex of degree $5$ can be realized
simply on six vertices.)
\end{exercise}

\begin{exercise}[Components by parity]
\label{ch01:ex-components}
Prove that in any graph the number of vertices of odd degree is even
\emph{within each connected component}, not merely globally. Deduce that a
connected graph with exactly two odd-degree vertices has an Euler trail between
them (you may cite the Euler-trail characterization informally; a full proof
comes later).
\end{exercise}

\begin{exercise}[Representation surgery]
\label{ch01:ex-transpose}
Given a digraph as an adjacency list, give an $\Oh(n+m)$ algorithm that outputs
the adjacency list of the \emph{transpose} graph $\Gph^{\mathsf T}$ (every edge
reversed). Prove the bound. Then explain why doing the same from an adjacency
matrix is $\Th(n^2)$ and from an edge list is $\Th(n+m)$.
\end{exercise}

\begin{exercise}[Implement CSR]
\label{ch01:ex-csr}
Implement \textsc{BuildCSR} (Algorithm~\ref{ch01:alg-csr}) and test it on the
undirected running example by first expanding each undirected edge into two
directed copies. Verify that $\mathrm{head}[n] = 2m$ and that scanning every
block reproduces all $2m$ directed edges exactly once.
\end{exercise}

\begin{exercise}[Break the twin trick]
\label{ch01:ex-twin}
Build an undirected graph as an edge array where the twin of edge $i$ is
$i \oplus 1$. Write a routine that, given an edge index, returns its twin, and a
routine that walks an Euler circuit using the twins. Then deliberately break the
layout (store one edge's copies at non-adjacent indices) and describe precisely
what goes wrong in the Euler walk. Why does the $i \oplus 1$ rule fail the moment
the pairing is not at $2i, 2i{+}1$?
\end{exercise}

\begin{exercise}[Bipartite test by parity]
\label{ch01:ex-bipartite}
Using the distance-parity coloring from the proof of
Theorem~\ref{ch01:bip-oddcycle}, write $\Oh(n+m)$ pseudocode that either outputs a
valid $2$-coloring or returns an odd cycle as a certificate of
non-bipartiteness. Argue that the returned cycle really is odd.
\end{exercise}

\begin{exercise}[Walks via matrix powers]
\label{ch01:ex-matrixpower}
Prove by induction that $(\Adjm^k)_{uv}$ equals the number of walks of length
exactly $k$ from $u$ to $v$. Use it to compute, for the running five-vertex
graph, the number of closed walks of length $3$ from vertex $3$ to itself, and
relate the trace $\sum_v (\Adjm^3)_{vv}$ to the number of triangles in the
graph (work out the exact constant of proportionality).
\end{exercise}

\begin{exercise}[Tree edge count, two ways]
\label{ch01:ex-treecount}
Give a second proof that a tree on $n$ vertices has $n-1$ edges, this time by
induction on $n$ using the existence of a leaf
(Corollary~\ref{ch01:forest-count}): remove a leaf and recurse. Where exactly
does the argument use acyclicity, and where does it use connectivity?
\end{exercise}

\begin{exercise}[Dense vs sparse crossover]
\label{ch01:ex-crossover}
At what density $m$ (as a function of $n$) do an adjacency matrix and an
adjacency list use asymptotically the same space? For a graph with $m = 4n$
edges and $n = 10^6$, by what factor does the matrix waste memory compared with
the list? Give actual numbers.
\end{exercise}

\begin{exercise}[Implicit graph blow-up]
\label{ch01:ex-implicit}
Consider the implicit graph whose vertices are binary strings of length $k$ and
where $\mathrm{succ}(s)$ flips any single bit (the hypercube $Q_k$). It has
$n = 2^k$ vertices and $m = k\,2^{k-1}$ edges. Following the successor rule
in Algorithm~\ref{ch01:alg-implicit-bfs}, trace the step-by-step distance
expansion from $0^k$ to $1^k$ on paper for $k = 3$, grouping vertices by their
Hamming distance from $0^k$, and confirm the shortest path has length $k$.
Then explain why storing this graph explicitly is impossible for $k = 40$,
and contrast the number of states touched to find a short path against storing
the full state space.
\end{exercise}
